\begin{table}
\caption{Supplementary Table S9. Sensitivity of the quality-performance associations to how the strata are cut. The primary analysis contrasts the within-source quartiles (<=Q25 against >=Q75) for the three continuous indicators; this repeats it at Q20 against Q80. The structural indicators are present-versus-absent in both and so are unchanged by the cut. The purpose is to show the principal associations are not an artefact of one arbitrary threshold.}
\begin{tabular}{lrrrrrrr}
\toprule
 & Q25/Q75 dAUPRC & Q25/Q75 lo & Q25/Q75 hi & sources & Q20/Q80 dAUPRC & Q20/Q80 lo & Q20/Q80 hi \\
indicator &  &  &  &  &  &  &  \\
\midrule
HF noise & 0.052300 & 0.041700 & 0.064100 & 4 & 0.056100 & 0.044700 & 0.068100 \\
baseline wander & 0.010300 & -0.001100 & 0.022300 & 4 & 0.009600 & -0.002000 & 0.022300 \\
mains & 0.056300 & 0.045200 & 0.066900 & 4 & 0.061600 & 0.049200 & 0.074000 \\
flat/clipped lead & -0.149700 & -0.190500 & -0.108300 & 1 & -0.149700 & -0.190500 & -0.108300 \\
QRS failure & -0.098500 & -0.129800 & -0.068200 & 2 & -0.098500 & -0.129800 & -0.068200 \\
RR implausibility & -0.126800 & -0.150600 & -0.102900 & 2 & -0.126800 & -0.150600 & -0.102900 \\
\bottomrule
\end{tabular}
\end{table}
